\section{Introducción}\label{introduction}

Estos son documentos de trabajo, no argumentos. Cada capítulo es una guía sobre un instrumento cartomántico concreto: su estructura, su vocabulario, su método. Nada de lo que aquí se expone necesita justificación, y no hace falta creer en nada de esto para poder utilizarlo.

La pregunta de por qué funciona la adivinación es razonable y tiene una respuesta práctica que no requiere ningún compromiso metafísico. Cada sistema de este libro ofrece un vocabulario simbólico estructurado ---un conjunto finito de fuerzas con nombre, organizadas en una arquitectura coherente--- y un método para aplicar ese vocabulario a una situación. El acto de aplicar el mapa no es pasivo. Requiere que el practicante seleccione, ordene e interprete las relaciones entre los símbolos, lo cual es otra forma de decir que exige que el practicante preste atención a la situación con una precisión que el pensamiento ordinario no impone. La estructura es el valor. Un buen sistema cartomántico no te dice qué pensar: te dice qué mirar y en qué relación.

Más allá de la explicación estructural, cada sistema de este libro tiene una lógica interna que amplifica, en lugar de sustituir, la percepción del practicante. El Tarot de la Golden Dawn ofrece un marco cabalístico lo suficientemente riguroso como para que el significado de la carta quede condicionado desde múltiples direcciones simultáneamente: Sephirah, elemento, decanato, letra hebrea y camino convergen todos en una sola fuerza, que no puede malinterpretarse fácilmente en ninguna de sus capas sin que las demás la corrijan. El Tarot Smith-Waite codifica el mismo marco bajo escenas pictóricas que se comunican directamente con la percepción, evitando la necesidad de calcular conscientemente las correspondencias. El Lenormand genera significado a través de la combinación: las cartas sueltas son incompletas y, por lo tanto, el lector siempre trabaja con relaciones en lugar de símbolos aislados, lo que refleja la estructura relacional de las situaciones reales. El sistema de naipes inglés es el más directo de los cuatro: cartas individuales en posiciones, leídas sin más, sin peso cosmológico.

Son cuatro instrumentos diferentes para preguntas diferentes, registros diferentes y profundidades diferentes. Una práctica completa utiliza la herramienta adecuada para la pregunta adecuada. El capítulo final de este libro aborda esa cuestión directamente.

Los capítulos siguen un orden descendente, de lo cosmológico a lo circunstancial: primero el Tarot de la Golden Dawn, segundo el Tarot de Smith-Waite, tercero el sistema de naipes inglés y cuarto el Oráculo Lenormand. No hay un orden obligatorio. Cada capítulo es independiente. La comparación corresponde al capítulo final y a la propia experiencia acumulada del practicante.

\begin{center}\rule{0.5\linewidth}{0.5pt}\end{center}
